\subsection{Deadlock Freedom}
\label{sec:deadlock-freedom-agda}

As we have seen in \Cref{sec:calculus}, the deadlock freedom property and
\Cref{thm:df} rest on some notions and predicates about processes which must be
formalised in Agda. First of all we need to define the notion of \emph{thread},
that is any process other than a cut. It is convenient to provide a more
fine-grained classification of threads, distinguishing between links and
input/output actions and sometimes also on whether such actions operate on free
or bound channels. We define predicates for each of these classes:

\begin{code}[hide]
{-# OPTIONS --rewriting #-}
open import Data.Unit using (tt)
open import Data.Sum
open import Data.Product using (_×_; _,_; ∃; ∃-syntax)
open import Data.List.Base using ([]; _∷_; [_])
open import Relation.Nullary using (¬_; contradiction)
open import Relation.Unary
open import Relation.Binary.PropositionalEquality using (refl)

open import Type
open import Context
open import Process
open import Reduction
open import Congruence
\end{code}

\begin{code}
data Link           : ∀{Γ} → Proc Γ → Set
data Input          : ∀{Γ} → Proc Γ → Set
data Output         : ∀{Γ} → Proc Γ → Set
data Delayed        : ∀{Γ} → Proc Γ → Set
data Server         : ∀{Γ} → Proc Γ → Set
data DelayedServer  : ∀{Γ} → Proc Γ → Set
\end{code}

The implementation of these predicates is not interesting since it is
essentially isomorphic to the relevant fragments of the \AgdaDatatype{Proc}
datatype. The only aspect that is worth pointing out here is that in
\AgdaDatatype{Input}, \AgdaDatatype{Output} and \AgdaDatatype{Server} the
channel being acted upon by the thread is the \emph{first} in $\Context$, hence
it is the \emph{most recently introduced} channel, whereas in
\AgdaDatatype{Delayed} and \AgdaDatatype{DelayedServer} the channel is not the
first. This allows us to distinguish those threads that, in the context of a
cut, operate on the channel bound by the cut or on a free channel. To clarify
this aspect, let us look at the implementation of the constructor
\AgdaInductiveConstructor{wait} in \AgdaDatatype{Input} and in
\AgdaDatatype{Delayed}. In the former predicate we have

\begin{AgdaAlign}
\begin{code}[hide]
data Link where
  link : ∀{Γ A} (p : Γ ≃ [ A ] + [ dual A ]) → Link (link (ch ⟨ p ⟩ ch))

data Input where
  fail : ∀{Γ Δ} (p : Γ ≃ [] + Δ) → Input (fail (ch ⟨ < p ⟩ tt))
\end{code}

\begin{code}
  wait : ∀{Γ Δ P} (p : Γ ≃ [] + Δ) → Input (wait (ch ⟨ < p ⟩ P))
\end{code}
where the use of the constructor \AgdaInductiveConstructor{<} indicates that the
thread operates on the most recent channel. In the latter predicate we have

\begin{code}[hide]
  case : ∀{Γ Δ A B P Q} (p : Γ ≃ [] + Δ) → Input (case {A} {B} (ch ⟨ < p ⟩ (P , Q)))
  join : ∀{Γ Δ A B P} (p : Γ ≃ [] + Δ) → Input (join {A} {B} (ch ⟨ < p ⟩ P))
  all  : ∀{A Γ Δ F} (p : Γ ≃ [] + Δ) → Input (all {A} (ch ⟨ < p ⟩ F))

data Output where
  close    : Output (close ch)
  select-l : ∀{Γ Δ A B P} (p : Γ ≃ [] + Δ) → Output (select {A} {B} (ch ⟨ < p ⟩ inj₁ P))
  select-r : ∀{Γ Δ A B P} (p : Γ ≃ [] + Δ) → Output (select {A} {B} (ch ⟨ < p ⟩ inj₂ P))
  fork     : ∀{Γ Δ Δ₁ Δ₂ A B P Q} (p : Γ ≃ [] + Δ) (q : Δ ≃ Δ₁ + Δ₂) → Output (fork {A} {B} (ch ⟨ < p ⟩ (P ⟨ q ⟩ Q)))
  ex       : ∀{A B Γ Δ P} (p : Γ ≃ [] + Δ) → Output (ex {A} {B} (ch ⟨ < p ⟩ P))
  client   : ∀{Γ Δ A P}   (p : Γ ≃ [] + Δ) → Output (client {A} (ch ⟨ < p ⟩ P))
  weaken   : ∀{Γ Δ A P}   (p : Γ ≃ [] + Δ) → Output (weaken {A} (ch ⟨ < p ⟩ P))
  contract : ∀{Γ Δ A P}   (p : Γ ≃ [] + Δ) → Output (contract {A} (ch ⟨ < p ⟩ P))

data Delayed where
  fail     : ∀{C Γ Δ} (p : Γ ≃ [ ⊤ ] + Δ) → Delayed (fail (ch ⟨ >_ {C} p ⟩ tt))
\end{code}

\begin{code}
  wait     : ∀{C Γ Δ P} (p : Γ ≃ [ ⊥ ] + Δ) → Delayed (wait (ch ⟨ >_ {C} p ⟩ P))
\end{code}
where the use of the constructor \AgdaInductiveConstructor{>} indicates that the
thread operates on a channel other than the most recent one. Note that here we
have to specify the type $C$ in front of $\Context$ or else Agda is unable to
automatically resolve some metavariables.

\begin{code}[hide]
  case     : ∀{Γ Δ C A B P} (p : Γ ≃ [ A & B ] + Δ) → Delayed (case {A} {B} (ch ⟨ >_ {C} p ⟩ P))
  select-l : ∀{Γ Δ C A B P} (p : Γ ≃ [ A ⊕ B ] + Δ) → Delayed (select (ch ⟨ >_ {C} p ⟩ inj₁ P))
  select-r : ∀{Γ Δ C A B P} (p : Γ ≃ [ A ⊕ B ] + Δ) → Delayed (select (ch ⟨ >_ {C} p ⟩ inj₂ P))
  join     : ∀{Γ Δ C A B P} (p : Γ ≃ [ A ⅋ B ] + Δ) → Delayed (join (ch ⟨ >_ {C} p ⟩ P))
  fork-l   : ∀{Γ Δ Δ₁ Δ₂ C A B P Q} (p : Γ ≃ [ A ⊗ B ] + Δ) (q : Δ ≃ Δ₁ + Δ₂) →
             Delayed (fork (ch ⟨ >_ {C} p ⟩ (P ⟨ < q ⟩ Q)))
  fork-r   : ∀{Γ Δ Δ₁ Δ₂ C A B P Q} (p : Γ ≃ [ A ⊗ B ] + Δ) (q : Δ ≃ Δ₁ + Δ₂) →
             Delayed (fork (ch ⟨ >_ {C} p ⟩ (P ⟨ > q ⟩ Q)))
  all      : ∀{A C Γ Δ F} (p : Γ ≃ [ `∀ A ] + Δ) → Delayed (all (ch ⟨ >_ {C} p ⟩ F))
  ex       : ∀{A B C Γ Δ P} (p : Γ ≃ [ `∃ A ] + Δ) → Delayed (ex {A} {B} (ch ⟨ >_ {C} p ⟩ P))
  client   : ∀{Γ Δ A C P} (p : Γ ≃ [ `? A ] + Δ) → Delayed (client (ch ⟨ >_ {C} p ⟩ P))
  weaken   : ∀{Γ Δ A C P} (p : Γ ≃ [ `? A ] + Δ) → Delayed (weaken (ch ⟨ >_ {C} p ⟩ P))
  contract : ∀{Γ Δ A C P} (p : Γ ≃ [ `? A ] + Δ) → Delayed (contract (ch ⟨ >_ {C} p ⟩ P))

data Server where
  server : ∀{Γ Δ A P} (p : Γ ≃ [] + Δ) (un : Un Δ) → Server (server {A} (ch ⟨ < p ⟩ (un , P)))

data DelayedServer where
  server : ∀{Γ Δ A C P} (p : Γ ≃ [ `! A ] + Δ) (un : Un Δ) →
           DelayedServer (server {A} (ch ⟨ >_ p ⟩ (un-∷ {C} un , P)))
\end{code}
\end{AgdaAlign}

The predicate \AgdaDatatype{Thread} is simply the disjoint union of all the
previous ones.

\begin{code}
data Thread {Γ} (P : Proc Γ) : Set where
  link     : Link P → Thread P
  delayed  : Delayed P → Thread P
  output   : Output P → Thread P
  input    : Input P → Thread P
  server   : Server P → Thread P
  dserver  : DelayedServer P → Thread P
\end{code}

Observability, reducibility and aliveness are defined in the expected way:

\begin{code}
Observable : ∀{Γ} → Proc Γ → Set
Observable P = ∃[ Q ] P ⊒ Q × Thread Q

Reducible : ∀{Γ} → Proc Γ → Set
Reducible P = ∃[ Q ] P ↝ Q

Alive : ∀{Γ} → Proc Γ → Set
Alive P = Observable P ⊎ Reducible P
\end{code}

\begin{code}[hide]
fail→thread : ∀{Γ Δ} (p : Γ ≃ [ ⊤ ] + Δ) → Thread (fail (ch ⟨ p ⟩ tt))
fail→thread (< p) = input (fail p)
fail→thread (> p) = delayed (fail p)

wait→thread : ∀{Γ Δ P} (p : Γ ≃ [ ⊥ ] + Δ) → Thread (wait (ch ⟨ p ⟩ P))
wait→thread (< p) = input (wait p)
wait→thread (> p) = delayed (wait p)

case→thread : ∀{A B Γ Δ P} (p : Γ ≃ [ A & B ] + Δ) → Thread (case (ch ⟨ p ⟩ P))
case→thread (< p) = input (case p)
case→thread (> p) = delayed (case p)

left→thread : ∀{A B Γ Δ P} (p : Γ ≃ [ A ⊕ B ] + Δ) → Thread (select (ch ⟨ p ⟩ inj₁ P))
left→thread (< p) = output (select-l p)
left→thread (> p) = delayed (select-l p)

right→thread : ∀{A B Γ Δ P} (p : Γ ≃ [ A ⊕ B ] + Δ) → Thread (select (ch ⟨ p ⟩ inj₂ P))
right→thread (< p) = output (select-r p)
right→thread (> p) = delayed (select-r p)

join→thread : ∀{A B Γ Δ P} (p : Γ ≃ [ A ⅋ B ] + Δ) → Thread (join (ch ⟨ p ⟩ P))
join→thread (< p) = input (join p)
join→thread (> p) = delayed (join p)

fork→thread : ∀{A B Γ Δ Δ₁ Δ₂ P Q} (p : Γ ≃ [ A ⊗ B ] + Δ) (q : Δ ≃ Δ₁ + Δ₂) → Thread (fork (ch ⟨ p ⟩ (P ⟨ q ⟩ Q)))
fork→thread (< p) q = output (fork p q)
fork→thread (> p) (< q) = delayed (fork-l p q)
fork→thread (> p) (> q) = delayed (fork-r p q)

all→thread : ∀{A Γ Δ F} (p : Γ ≃ [ `∀ A ] + Δ) → Thread (all (ch ⟨ p ⟩ F))
all→thread (< p) = input (all p)
all→thread (> p) = delayed (all p)

ex→thread : ∀{A B Γ Δ P} (p : Γ ≃ [ `∃ A ] + Δ) → Thread (ex {A} {B} (ch ⟨ p ⟩ P))
ex→thread (< p) = output (ex p)
ex→thread (> p) = delayed (ex p)

server→thread : ∀{A Γ Δ P} (p : Γ ≃ [ `! A ] + Δ) (un : Un Δ) → Thread (server (ch ⟨ p ⟩ (un , P)))
server→thread (< p) un = server (server p un)
server→thread (> p) (un-∷ un) = dserver (server p un)

client→thread : ∀{A Γ Δ P} (p : Γ ≃ [ `? A ] + Δ) → Thread (client (ch ⟨ p ⟩ P))
client→thread (< p) = output (client p)
client→thread (> p) = delayed (client p)

weaken→thread : ∀{A Γ Δ P} (p : Γ ≃ [ `? A ] + Δ) → Thread (weaken {A = A} (ch ⟨ p ⟩ P))
weaken→thread (< p) = output (weaken p)
weaken→thread (> p) = delayed (weaken p)

contract→thread : ∀{A Γ Δ P} (p : Γ ≃ [ `? A ] + Δ) → Thread (contract (ch ⟨ p ⟩ P))
contract→thread (< p) = output (contract p)
contract→thread (> p) = delayed (contract p)
\end{code}

In order to prove that every (well-typed) process is alive, it is convenient to
define a ``canonical'' form for cuts, that is a form that matches at least one
of the l.h.s of one of the rules for structural pre-congruence or reduction in
\Cref{tab:semantics}. This is the point where the fine-grained classification of
threads introduced earlier \REVISION{pays off: instead of considering all the 21
forms of cuts that can be rewritten/reduced according to the rules in
\Cref{tab:semantics}, we only need to consider four major classes of cuts:}

\begin{code}
data CanonicalCut {Γ} : Proc Γ → Set where
  cc-link     : ∀{Γ₁ Γ₂ A P Q} (p : Γ ≃ Γ₁ + Γ₂) →
                Link P → CanonicalCut (cut {A} (P ⟨ p ⟩ Q))
  cc-redex    : ∀{Γ₁ Γ₂ A P Q} (p : Γ ≃ Γ₁ + Γ₂) →
                Output P → (Input ∪ Server) Q →
                CanonicalCut (cut {A} (P ⟨ p ⟩ Q))
  cc-delayed  : ∀{Γ₁ Γ₂ A P Q} (p : Γ ≃ Γ₁ + Γ₂) →
                Delayed P → CanonicalCut (cut {A} (P ⟨ p ⟩ Q))
  cc-servers  : ∀{Γ₁ Γ₂ A P Q} (p : Γ ≃ Γ₁ + Γ₂) →
                DelayedServer P → Server Q →
                CanonicalCut (cut {A} (P ⟨ p ⟩ Q))
\end{code}

A \emph{canonical cut} $\Cut[A]\x{P}{Q}$ has one of these forms:
\begin{itemize}
\item $P$ is a link (\AgdaInductiveConstructor{cc-link}), or
\item $P$ performs an output on $x$ and $Q$ performs an input on $x$
  (\AgdaInductiveConstructor{cc-redex}), or
\item $P$ operates on a channel other than $x$ and is not a server
(\AgdaInductiveConstructor{cc-delayed}), or
\item both $P$ and $Q$ are servers and $P$ operates on a channel other than $x$
(\AgdaInductiveConstructor{cc-servers}).
\end{itemize}

\REVISION{Note that we} have to distinguish servers from the other input
operations because the structural pre-congruence rule \SServer can only be
applied when the two sub-processes of a cut are both servers.

\begin{code}[hide]
output-output : ∀{A Γ Δ P Q} → ¬ (Output {A ∷ Γ} P × Output {dual A ∷ Δ} Q)
output-output (close , ())

output-delayed-server : ∀{A Γ Δ P Q} → ¬ (Output {A ∷ Γ} P × DelayedServer {dual A ∷ Δ} Q)
output-delayed-server (close , ())

input-input : ∀{A Γ Δ P Q} → ¬ (Input {A ∷ Γ} P × Input {dual A ∷ Δ} Q)
input-input (fail _ , ())

input-server : ∀{A Γ Δ P Q} → ¬ (Input {A ∷ Γ} P × Server {dual A ∷ Δ} Q)
input-server (fail _ , ())

input-delayed-server : ∀{A Γ Δ P Q} → ¬ (Input {A ∷ Γ} P × DelayedServer {dual A ∷ Δ} Q)
input-delayed-server (fail _ , ())

server-server : ∀{A Γ Δ P Q} → ¬ (Server {A ∷ Γ} P × Server {dual A ∷ Δ} Q)
server-server (server _ _ , ())

delayed-server-delayed-served : ∀{A Γ Δ P Q} → ¬ (DelayedServer {A ∷ Γ} P × DelayedServer {dual A ∷ Δ} Q)
delayed-server-delayed-served (server _ _ , ())
\end{code}

Every cut $\Cut[A]\x{P}{Q}$ where both $P$ and $Q$ are threads can be rewritten
into a canonical cut using structural pre-congruence:

\begin{code}
canonical-cut : ∀{A Γ Γ₁ Γ₂ P Q} (p : Γ ≃ Γ₁ + Γ₂) →
                Thread P → Thread Q →
                ∃[ R ] CanonicalCut R × cut {A} (P ⟨ p ⟩ Q) ⊒ R
\end{code}

\begin{code}[hide]
canonical-cut pc (link x) Qt = _ , cc-link pc x , s-refl
canonical-cut pc Pt (link y) = _ , cc-link (+-comm pc) y , s-comm pc
canonical-cut pc (delayed x) Qt = _ , cc-delayed pc x , s-refl
canonical-cut pc Pt (delayed y) = _ , cc-delayed (+-comm pc) y , s-comm pc
canonical-cut pc (output x) (output y) = contradiction (x , y) output-output
canonical-cut pc (output x) (input y) = _ , cc-redex pc x (inj₁ y) , s-refl
canonical-cut pc (output x) (server y) = _ , cc-redex pc x (inj₂ y) , s-refl
canonical-cut pc (output x) (dserver y) = contradiction (x , y) output-delayed-server
canonical-cut pc (input x) (output y) = _ , cc-redex (+-comm pc) y (inj₁ x) , s-comm pc
canonical-cut pc (input x) (input y) = contradiction (x , y) input-input
canonical-cut pc (input x) (server y) = contradiction (x , y) input-server
canonical-cut pc (input x) (dserver y) = contradiction (x , y) input-delayed-server
canonical-cut pc (server x) (output y) = _ , cc-redex (+-comm pc) y (inj₂ x) , s-comm pc
canonical-cut pc (server x) (input y) = contradiction (y , x) input-server
canonical-cut pc (server x) (server y) = contradiction (x , y) server-server
canonical-cut pc (server x) (dserver y) = _ , cc-servers (+-comm pc) y x , s-comm pc
canonical-cut pc (dserver x) (output y) = contradiction (y , x) output-delayed-server
canonical-cut pc (dserver x) (input y) = contradiction (y , x) input-delayed-server
canonical-cut pc (dserver x) (server y) = _ , cc-servers pc x y , s-refl
canonical-cut pc (dserver x) (dserver y) = contradiction (x , y) delayed-server-delayed-served

⊒Alive : ∀{Γ} {P Q : Proc Γ} → P ⊒ Q → Alive Q → Alive P
⊒Alive pcong (inj₁ (_ , x , th)) = inj₁ (_ , s-tran pcong x , th)
⊒Alive pcong (inj₂ (_ , red)) = inj₂ (_ , r-cong pcong red)
\end{code}

It is easy to prove that every canonical cut is alive, either reducing it or
applying structural pre-congruence to rewrite it into a thread.

\begin{code}
canonical-cut-alive : ∀{Γ} {C : Proc Γ} → CanonicalCut C → Alive C
\end{code}

\begin{code}[hide]
canonical-cut-alive (cc-link pc (link (< > •))) = inj₂ (_ , r-link pc)
canonical-cut-alive (cc-link pc (link (> < •))) =
  inj₂ (_ , r-cong (s-cong pc (s-link _) s-refl) (r-link pc))
canonical-cut-alive (cc-redex pc close (inj₁ (wait p))) with +-empty-l p | +-empty-l pc
... | refl | refl = inj₂ (_ , r-close pc p)
canonical-cut-alive (cc-redex pc (select-l p) (inj₁ (case q))) with +-empty-l p | +-empty-l q
... | refl | refl = inj₂ (_ , r-select-l pc p q)
canonical-cut-alive (cc-redex pc (select-r p) (inj₁ (case q))) with +-empty-l p | +-empty-l q
... | refl | refl = inj₂ (_ , r-select-r pc p q)
canonical-cut-alive (cc-redex pc (fork p q) (inj₁ (join r))) with +-empty-l p | +-empty-l r
... | refl | refl = inj₂ (_ , r-fork pc r q p)
canonical-cut-alive (cc-redex pc (ex p) (inj₁ (all q))) with +-empty-l p | +-empty-l q
... | refl | refl = inj₂ (_ , r-exists pc p q)
canonical-cut-alive (cc-redex pc (client p) (inj₂ (server q un))) with +-empty-l p | +-empty-l q
... | refl | refl = inj₂ (_ , r-client pc p q un)
canonical-cut-alive (cc-redex pc (weaken p) (inj₂ (server q un))) with +-empty-l p | +-empty-l q
... | refl | refl = inj₂ (_ , r-weaken pc p q un)
canonical-cut-alive (cc-redex pc (contract p) (inj₂ (server q un))) with +-empty-l p | +-empty-l q
... | refl | refl = inj₂ (_ , r-contract pc p q un)
canonical-cut-alive (cc-delayed pc (fail p)) =
  let _ , _ , p' = +-assoc-l pc p in
  inj₁ (_ , s-fail pc p , fail→thread p')
canonical-cut-alive (cc-delayed pc (wait p)) =
  let _ , _ , p' = +-assoc-l pc p in
  inj₁ (_ , s-wait pc p , wait→thread p')
canonical-cut-alive (cc-delayed pc (case p)) =
  let _ , _ , p' = +-assoc-l pc p in
  inj₁ (_ , s-case pc p , case→thread p')
canonical-cut-alive (cc-delayed pc (join p)) =
  let _ , _ , p' = +-assoc-l pc p in
  inj₁ (_ , s-join pc p , join→thread p')
canonical-cut-alive (cc-delayed pc (select-l p)) =
  let _ , _ , p' = +-assoc-l pc p in
  inj₁ (_ , s-select-l pc p , left→thread p')
canonical-cut-alive (cc-delayed pc (select-r p)) =
  let _ , _ , p' = +-assoc-l pc p in
  inj₁ (_ , s-select-r pc p , right→thread p')
canonical-cut-alive (cc-delayed p (fork-l q r)) =
  let _ , p' , q' = +-assoc-l p q in
  let _ , p'' , r' = +-assoc-l p' r in
  let _ , q'' , r'' = +-assoc-r r' (+-comm p'') in
  inj₁ (_ , s-fork-l p q r , fork→thread q' r'')
canonical-cut-alive (cc-delayed p (fork-r q r)) =
  let _ , p' , q' = +-assoc-l p q in
  let _ , p'' , r' = +-assoc-l p' r in
  inj₁ (_ , s-fork-r p q r , fork→thread q' r')
canonical-cut-alive (cc-delayed pc (all p)) =
  let _ , _ , p' = +-assoc-l pc p in
  inj₁ (_ , s-all pc p , all→thread p')
canonical-cut-alive (cc-delayed pc (ex p)) =
  let _ , _ , p' = +-assoc-l pc p in
  inj₁ (_ , s-ex pc p , ex→thread p')
canonical-cut-alive (cc-delayed pc (client p)) =
  let _ , _ , p' = +-assoc-l pc p in
  inj₁ (_ , s-client pc p , client→thread p')
canonical-cut-alive (cc-delayed pc (weaken p)) =
  let _ , _ , p' = +-assoc-l pc p in
  inj₁ (_ , s-weaken pc p , weaken→thread p')
canonical-cut-alive (cc-delayed pc (contract p)) =
  let _ , _ , p' = +-assoc-l pc p in
  inj₁ (_ , s-contract pc p , contract→thread p')
canonical-cut-alive (cc-servers pc (server p un) (server q un')) with +-empty-l q
... | refl =
  let _ , pc' , p' = +-assoc-l pc p in
  inj₁ (_ , s-server pc p q un un' , server→thread p' (∗-un (un ⟨ pc' ⟩ un')))
\end{code}

Now deadlock freedom for $P$ can be proved by induction on $P$.

\begin{code}
deadlock-freedom : ∀{Γ} (P : Proc Γ) → Alive P
\end{code}

When $P$ is a thread, then it is obviously observable and hence alive. When $P$
is a cut $\Cut[A]\x{P}{Q}$, \AgdaFunction{deadlock-freedom} is applied
recursively to $P$ and to $Q$, in turn. If either of these applications yields a
reduction, then the whole cut is reducible and therefore alive. If both
applications yield a thread, then we conclude that the cut is alive first
rewriting it into a canonical cut with \AgdaFunction{canonical-cut} and then
using \AgdaFunction{canonical-cut-alive}.

\begin{code}[hide]
deadlock-freedom (link (ch ⟨ p ⟩ ch)) = inj₁ (_ , s-refl , link (link p))
deadlock-freedom (fail (ch ⟨ p ⟩ _)) = inj₁ (_ , s-refl , fail→thread p)
deadlock-freedom (wait (ch ⟨ p ⟩ _)) = inj₁ (_ , s-refl , wait→thread p)
deadlock-freedom (close ch) = inj₁ (_ , s-refl , output close)
deadlock-freedom (case (ch ⟨ p ⟩ _)) = inj₁ (_ , s-refl , case→thread p)
deadlock-freedom (select (ch ⟨ p ⟩ inj₁ _)) = inj₁ (_ , s-refl , left→thread p)
deadlock-freedom (select (ch ⟨ p ⟩ inj₂ _)) = inj₁ (_ , s-refl , right→thread p)
deadlock-freedom (join (ch ⟨ p ⟩ _)) = inj₁ (_ , s-refl , join→thread p)
deadlock-freedom (fork (ch ⟨ p ⟩ (P ⟨ q ⟩ Q))) = inj₁ (_ , s-refl , fork→thread p q)
deadlock-freedom (all (ch ⟨ p ⟩ F)) = inj₁ (_ , s-refl , all→thread p)
deadlock-freedom (ex (ch ⟨ p ⟩ P)) = inj₁ (_ , s-refl , ex→thread p)
deadlock-freedom (server (ch ⟨ p ⟩ (un , _))) = inj₁ (_ , s-refl , server→thread p un)
deadlock-freedom (client (ch ⟨ p ⟩ _)) = inj₁ (_ , s-refl , client→thread p)
deadlock-freedom (weaken (ch ⟨ p ⟩ _)) = inj₁ (_ , s-refl , weaken→thread p)
deadlock-freedom (contract (ch ⟨ p ⟩ _)) = inj₁ (_ , s-refl , contract→thread p)
deadlock-freedom (cut (P ⟨ p ⟩ Q)) with deadlock-freedom P
... | inj₂ (_ , red) = inj₂ (_ , r-cut p red)
... | inj₁ (_ , Pc , Pt) with deadlock-freedom Q
... | inj₂ (_ , red) = inj₂ (_ , r-cong (s-comm p) (r-cut (+-comm p) red))
... | inj₁ (_ , Qc , Qt) with canonical-cut p Pt Qt
... | _ , cc , pcong = ⊒Alive (s-tran (s-cong p Pc Qc) pcong) (canonical-cut-alive cc)
\end{code}
